\documentclass[10pt,a4paper]{article}

\usepackage[margin=1.2cm]{geometry}
\usepackage[T1]{fontenc}
\usepackage[utf8]{inputenc}
\usepackage{lmodern}
\usepackage{microtype}
\usepackage{enumitem}
\usepackage{xcolor}
\usepackage{graphicx}
\usepackage{parskip}

\setlength{\parindent}{0pt}
\setlength{\parskip}{0.28em}
\setlist[itemize]{leftmargin=*, itemsep=0pt, topsep=1pt}

\begin{document}

\begin{center}
{\LARGE \textbf{SHA-256 Hardware Acceleration in Croc}}\\[2pt]
{\large Project Outline for VLSI 2 FS2026}\\[2pt]
\textbf{Radu-Constantin Cretu, Alexander Huwiler}
\end{center}

\vspace{0.15em}
\small
\textbf{Project objective.}
We propose a Croc-based SoC improvement that adds a dedicated SHA-256 hardware
accelerator to the \texttt{user\_domain}. In the baseline exercise design,
SHA-256 would have to run entirely in software on the RISC-V core. Our project
offloads the SHA-256 compression step to custom hardware, turning Croc into a
small hardware/software co-designed platform with a measurable application-level
improvement. The proposed extension is therefore a genuine SoC-level addition:
it introduces a new user-domain peripheral, a new software-visible interface,
and a concrete workload for which the impact of the improvement can be measured.

\textbf{Design improvement.}
The accelerator is exposed as a memory-mapped peripheral at
\texttt{0x2000\_1000}. Software writes a padded 512-bit message block into 16
input registers, triggers the computation through a control register, and reads
the resulting 256-bit digest back from 8 output registers. The current design
supports both single-block hashing and chained multi-block hashing, so longer
messages can be processed block by block. This keeps the interface simple while
still providing a meaningful architectural extension over the base Croc design.
In the current version, software performs the message padding while the hardware
implements the SHA-256 compression and state chaining. This choice keeps the
accelerator compact enough for backend implementation while still moving the
performance-critical part of the algorithm into hardware. In addition, the
required user ROM at \texttt{0x2000\_0000} has been added so that the chip can
return the student names in the format expected by the course infrastructure.

\textbf{Relevance and originality.}
SHA-256 is a realistic accelerator target because it is computationally heavy,
easy to verify against standard references, and structurally well suited for a
custom datapath. The project therefore aligns well with the course goals:
functionality can be demonstrated unambiguously, performance can be compared
against a software baseline, and the final design remains suitable for the full
implementation flow from RTL to physical design. The proposed accelerator is not
just an isolated module, but a complete Croc extension involving user-domain
integration, software support, verification, and later backend tradeoffs in area
and timing.

\textbf{System-level structure.}
At the SoC level, the project adds a new accelerator block in the
\texttt{user\_domain} address space while leaving the existing Croc core and
baseline peripherals intact. Software communicates with the accelerator through
a simple programming model based on control, status, block-input, and digest
registers. This makes the extension easy to exercise from bare-metal software,
but still representative of a realistic accelerator integration problem:
address decoding, register-map design, interrupt support, and interaction
between the processor and a dedicated datapath. In this sense, the project is
not only about implementing SHA-256 itself, but about integrating a useful
accelerator into a complete SoC in a way that remains compatible with the
standard Croc flow.

\textbf{Implementation status.}
A first working prototype is already integrated into Croc. The RTL, MMIO
interface, driver, and bare-metal software tests are in place. Functional
verification in Verilator already covers standard single-block and multi-block
SHA-256 vectors, interrupt behavior, and relevant block-boundary cases. A first
benchmark shows a hardware speedup of about \texttt{14.4x} over the software-only
implementation for the current one-block \texttt{"abc"} case. Synthesis also runs
cleanly, which gives us a solid frontend checkpoint before focusing on backend.
This means that the project is already beyond the idea stage: the accelerator is
integrated into the Croc environment, usable from software, and supported by a
working verification and synthesis flow.

\textbf{Evaluation plan.}
The final project will be evaluated along the following axes:
\begin{itemize}
  \item \textbf{Functional correctness:} compare the accelerator against
  standard SHA-256 test vectors and software reference behavior.
  \item \textbf{Performance improvement:} measure speedup in cycle count
  relative to the software-only implementation.
  \item \textbf{Implementation cost:} collect synthesis and backend results,
  with particular attention to area and timing impact.
\end{itemize}

\textbf{Expected outcome.}
The intended final result is a Croc variant that demonstrates a clear,
documented improvement over the exercise baseline. On the functional side, the
accelerator should execute the standard SHA-256 compression correctly for
single-block and multi-block workloads. On the performance side, it should show
a measurable speedup over the software-only reference running on the Croc CPU.
On the implementation side, the design should remain compatible with the fixed
project constraints and be pushed as far as possible through the backend flow,
so that the final report can discuss the accelerator not only as an algorithmic
improvement, but also as a physically implementable SoC extension.

\textbf{Next steps.}
The next phase is to adapt the backend flow to the fixed 2026 chip-size and pad
constraints, run the OpenROAD-based physical implementation flow, and collect
area and timing results. The final report will then document the architecture,
RTL modifications, verification approach, measured speedup, and the main backend
challenges encountered during implementation. Concretely, the remaining work is
to move from a verified frontend implementation to a backend-complete project:
we must align the design with the fixed physical constraints of the 2026
exercise framework, push the design through place-and-route, and assess the
implementation cost of the improvement in silicon terms. This will allow the
final project to argue not only that the accelerator works, but also that it is
a technically meaningful and well-justified addition to the Croc SoC.

\end{document}
